\documentclass[12pt]{book}
\usepackage{amsmath}
\usepackage{amssymb}
\usepackage{geometry}
\geometry{margin=1in}

\title{Volume 04: Electromagnetism and Magnetic Phenomena \\ Book 02: Magnetic Moments and Helical Wakes}
\author{James C. Harvey}
\date{December 2025}

\begin{document}
\maketitle
\tableofcontents

\chapter{Magnetism as Helical Wake}

\section*{Abstract}
Magnetism in SDT is the helical wake of circulation in the matrix. Magnetic moments arise from toroidal
circulation and are proportional to curvature, circulation, and coupling efficiency. This book formalizes
the magnetic moment expression and its sign.

\section{Helical Wake Definition}
The magnetic field is the helical wake pattern produced by circulation. Its handedness is the sign of
the magnetic moment.

\section{Magnetic Moment from Circulation}
The moment scales with geometric factors:
\begin{equation}
\mu \propto \Gamma \kappa (1-\eta).
\end{equation}
The sign is determined by circulation chirality.

\section{Negative Moments}
A negative magnetic moment indicates reversed circulation relative to the chosen axis. This is a direct
geometric signature, not a statistical artifact.

\chapter{Magnetic Field Geometry}

\section*{Abstract}
The magnetic field is a spatial map of helical wake density. It is the directional curvature of matrix
flux around circulating boundaries.

\section{Field from Circulation}
For a circulating boundary, the wake amplitude scales as
\begin{equation}
B \propto \Gamma \kappa \frac{(1-\eta)}{r^2}.
\end{equation}
This is the far-field decay of helical wake intensity.

\section{Scaling Ratio}
For two radii $r_1$ and $r_2$ with identical circulation and curvature,
\begin{equation}
\frac{B_1}{B_2} = \left(\frac{r_2}{r_1}\right)^2,
\end{equation}
which encodes the inverse-square decay in SDT geometry.

\section{Dipole Form}
At distances large compared to boundary scale, the field is dipolar:
\begin{equation}
\mathbf{B}(\mathbf{r}) = \frac{\mu_0}{4\pi}\frac{3(\boldsymbol{\mu}\cdot \hat{\mathbf{r}})\hat{\mathbf{r}} - \boldsymbol{\mu}}{r^3},
\end{equation}
with $\mu_0$ interpreted as a matrix compliance constant.

\chapter{Induction as Wake Reconfiguration}

\section*{Abstract}
Induction is wake reconfiguration in the matrix. Changing circulation modifies local flux topology and
drives opposing responses.

\section{Induction Law}
\begin{equation}
\nabla \times \mathbf{E} = -\frac{\partial \mathbf{B}}{\partial t}.
\end{equation}
This relation is a geometric necessity of wake continuity.

\chapter{Magnetic Moments from Geometry}

\section*{Abstract}
Magnetic moments are not empirical constants in SDT; they are computable from geometry. This chapter
summarizes the geometric formula and points to the benchmark derivations.

\section{Moment Expression}
\begin{equation}
\mu = C_\mu \Gamma \kappa (1-\eta),
\end{equation}
where $C_\mu$ is the geometric conversion constant defined by matrix scaling.

\section{Numerical Template}
The moment can be evaluated numerically by substituting geometric parameters:
\begin{equation}
\mu = C_\mu \Gamma \kappa (1-\eta)\ \text{J/T}.
\end{equation}
All benchmark computations apply this template without fitted parameters.

\section{Canonical Derivations}
Full particle derivations are provided in \texttt{SDT/benchmarks/Grok\_Benchmarks/Magnetic\_Moments\_SDT\_Real\_Derivation.md}
and the associated calculation scripts.
\end{document}
